\section{附录：数据采集与图像处理}
\subsection{视频下载与帧提取}
为了获得清晰的封面与 key frames，我们通过 \texttt{yt-dlp -f 398+251 -o lecture11.webm} 下载精度 720p 的原始视频（视频链接 \texttt{https://www.youtube.com/watch?v=\_MNlLhU33H0}）；然后在 00:02:35、00:36:00、00:44:40 与 01:01:30 四个时间点调用 \texttt{ffmpeg -ss ... -frames:v 1 -update 1}，依次生成 frame-01–frame-04。如下表列出关键命令与用途。

\begin{table}[H]
\centering
\begin{tabular}{p{6cm}p{9cm}}
\toprule
命令 & 说明 \\
\midrule
\texttt{yt-dlp -f 398+251 -o lecture11.webm https://www.youtube.com/watch?v=\_MNlLhU33H0} & 下载 720p+audio 视频，为后续截图与封面提供素材。 \\
\texttt{ffmpeg -ss 00:02:35 -i lecture11.webm -frames:v 1 -update 1 frame-01-selfreward.jpg} & 捕捉 self-rewarding 章节 key frame。 \\
\texttt{ffmpeg -ss 00:36:00 -i lecture11.webm -frames:v 1 -update 1 frame-02-iterative.jpg} & 捕捉 iterative pipeline 图示。 \\
\texttt{ffmpeg -ss 00:44:40 -i lecture11.webm -frames:v 1 -update 1 frame-03-selfinstr.jpg} & Self-Instruct 示例帧。 \\
\texttt{ffmpeg -ss 01:01:30 -i lecture11.webm -frames:v 1 -update 1 frame-04-meta.jpg} & Meta judge 画面。 \\
\bottomrule
\end{tabular}
\caption{关键采集命令与图像。}
\end{table}

\subsection{字幕清洗与提纲}
我们通过 Python 清洗 SRT，删去重复段落、保留时间戳，并按 2500 行/块提取主题句以辅助写作。脚本核心通过 \texttt{re.split(r'\\n\\n+')} 拆分区块、抽出第一行时间，并用 \texttt{prev} 去重保留精华。得到的 \texttt{subs\_clean.txt} 在笔记写作中作为 chronology 的来源 —— 它包含 \texttt{Self-Rewarding}、\texttt{Meta Judge} 等段落关键词，便于按教学逻辑分类。

\subsection{本章小结}
附录记录的下载、截图与字幕处理流程构成了后续写作的后方保障，确保每张图像与每段文字都可追溯到具体时间点和脚本。

